%======================================================================
%  plotstyle.tex
%  Unified figure and table styling for the Masterarbeit.
%
%  Everything visual — colours, markers, axis geometry, table rules,
%  caption formatting — is defined here and nowhere else. To change the
%  look of every figure or table in the thesis, edit this file only.
%
%  Requires (already loaded in head.tex): pgfplots, xcolor, booktabs,
%  caption, multirow, tabularx.
%  Input this file at the end of the preamble:  %======================================================================
%  plotstyle.tex
%  Unified figure and table styling for the Masterarbeit.
%
%  Everything visual — colours, markers, axis geometry, table rules,
%  caption formatting — is defined here and nowhere else. To change the
%  look of every figure or table in the thesis, edit this file only.
%
%  Requires (already loaded in head.tex): pgfplots, xcolor, booktabs,
%  caption, multirow, tabularx.
%  Input this file at the end of the preamble:  %======================================================================
%  plotstyle.tex
%  Unified figure and table styling for the Masterarbeit.
%
%  Everything visual — colours, markers, axis geometry, table rules,
%  caption formatting — is defined here and nowhere else. To change the
%  look of every figure or table in the thesis, edit this file only.
%
%  Requires (already loaded in head.tex): pgfplots, xcolor, booktabs,
%  caption, multirow, tabularx.
%  Input this file at the end of the preamble:  %======================================================================
%  plotstyle.tex
%  Unified figure and table styling for the Masterarbeit.
%
%  Everything visual — colours, markers, axis geometry, table rules,
%  caption formatting — is defined here and nowhere else. To change the
%  look of every figure or table in the thesis, edit this file only.
%
%  Requires (already loaded in head.tex): pgfplots, xcolor, booktabs,
%  caption, multirow, tabularx.
%  Input this file at the end of the preamble:  \input{plotstyle}
%======================================================================


%----------------------------------------------------------------------
% 1. COLOUR PALETTE
%----------------------------------------------------------------------
% Okabe-Ito palette: distinguishable under all common forms of colour
% vision deficiency. The L* value in each comment is the CIE lightness;
% series that must stay apart in greyscale printing are paired so that
% their L* values differ by at least 20.
%----------------------------------------------------------------------
\definecolor{tcBlue}{HTML}{0072B2}        % L* 47  — primary series
\definecolor{tcVermillion}{HTML}{D55E00}  % L* 52  — secondary series
\definecolor{tcTeal}{HTML}{009E73}        % L* 58  — third series
\definecolor{tcAmber}{HTML}{E69F00}       % L* 71  — second bar series
\definecolor{tcSky}{HTML}{56B4E9}         % L* 72  — spare
\definecolor{tcPurple}{HTML}{CC79A7}      % L* 61  — spare
\definecolor{tcGrey}{HTML}{595959}        % L* 38  — literature / prior work
\definecolor{tcGridGrey}{HTML}{BFBFBF}    % grid lines
\definecolor{tcRule}{HTML}{404040}        % axis lines and table rules

% Institutional colour, kept for the title page and header only.
% (HTWGreen is defined in style.tex and is deliberately not used as a
% data colour: it is hard to separate from teal in greyscale.)


%----------------------------------------------------------------------
% 2. SEMANTIC SERIES STYLES
%----------------------------------------------------------------------
% Colour carries MEANING and is constant across the whole thesis:
%
%   primary    blue        this work, main condition   (French, 0.5 mL,
%                                                       as-analyzed, pristine)
%   secondary  vermillion  this work, second condition (Algerian, 0.1 mL)
%   tertiary   teal        this work, third condition
%   emphasis   amber       derived / corrected quantity in BAR charts
%                          (ash-free, activated, sulfur-loaded)
%   reference  grey        literature or the preceding project work
%
% Shape and dash pattern duplicate that information so the figures stay
% readable in black-and-white print.
%----------------------------------------------------------------------
\pgfplotsset{
  % --- line series -----------------------------------------------------
  s primary/.style={
    tcBlue, thick, mark=*, mark size=2.3pt,
    mark options={fill=tcBlue, draw=tcBlue}},
  s secondary/.style={
    tcVermillion, thick, mark=square*, mark size=2.2pt,
    mark options={fill=tcVermillion, draw=tcVermillion}},
  s tertiary/.style={
    tcTeal, thick, mark=triangle*, mark size=2.8pt,
    mark options={fill=tcTeal, draw=tcTeal}},
  s quaternary/.style={
    tcAmber, thick, mark=diamond*, mark size=2.8pt,
    mark options={fill=tcAmber, draw=tcAmber}},
  s reference/.style={
    tcGrey, thick, densely dashed, mark=square*, mark size=2.2pt,
    mark options={solid, fill=tcGrey, draw=tcGrey}},
  % Open marker + dashed line = the same quantity on a different basis
  % (e.g. wet vs. moisture-corrected). Hue stays with the condition.
  s open/.style={
    dashed, mark options={solid, fill=white, line width=0.8pt}},
  %
  % --- bar series ------------------------------------------------------
  b primary/.style={fill=tcBlue,      draw=tcBlue!70!black,   line width=0.4pt},
  b emphasis/.style={fill=tcAmber,    draw=tcAmber!70!black,  line width=0.4pt},
  b secondary/.style={fill=tcVermillion, draw=tcVermillion!70!black, line width=0.4pt},
  b reference/.style={fill=tcGrey!30, draw=tcGrey,            line width=0.4pt},
}


%----------------------------------------------------------------------
% 3. AXIS STYLES
%----------------------------------------------------------------------
% One geometry for every plot in the thesis. The in-axis title is set once
% here, so every figure title is the same size and sits the same distance
% above the axis.
%----------------------------------------------------------------------
\pgfplotsset{
  thesisaxis/.style={
    width=0.86\textwidth,
    height=0.54\textwidth,
    title style={font=\normalsize, yshift=2pt},
    axis lines*=box,
    axis line style={tcRule, line width=0.5pt},
    tick align=outside,
    major tick length=3pt,
    minor tick num=0,
    tick style={tcRule, line width=0.4pt},
    ymajorgrids=true,
    grid style={tcGridGrey, line width=0.3pt},
    tick label style={font=\small},
    label style={font=\small},
    legend cell align=left,
    legend columns=2,
    legend style={
      font=\small,
      draw=tcGridGrey,
      fill=white,
      at={(0.5,-0.20)}, anchor=north,
      /tikz/every even column/.append style={column sep=0.6cm}},
    error bars/error bar style={tcRule, line width=0.6pt},
    error bars/error mark options={tcRule, rotate=90, mark size=2pt},
  },
  % line / scatter plots
  thesisline/.style={thesisaxis},
  % square-ish panel for scatter diagrams (e.g. van Krevelen)
  thesissquare/.style={
    thesisaxis,
    width=0.66\textwidth, height=0.54\textwidth,
    grid=both,
    minor grid style={tcGridGrey!50, line width=0.2pt}},
  % grouped bar charts
  thesisbar/.style={
    thesisaxis,
    ybar,
    area legend,
    bar width=18pt,
    ymin=0,
    enlarge x limits=0.28,
    xtick=data,
    tick pos=left,
    clip=false},
  % value labels on top of bars
  barlabels/.style={
    nodes near coords,
    nodes near coords align={vertical},
    every node near coord/.append style={
      font=\footnotesize, text=black,
      /pgf/number format/.cd, fixed, fixed zerofill, precision=1}},
  % overlay axis for a second y scale (no grid, no x axis)
  thesisaxis right/.style={
    width=0.86\textwidth,
    height=0.54\textwidth,
    axis y line*=right,
    axis x line=none,
    axis line style={tcRule, line width=0.5pt},
    tick align=outside,
    major tick length=3pt,
    tick style={tcRule, line width=0.4pt},
    tick label style={font=\small},
    label style={font=\small},
    error bars/error bar style={tcRule, line width=0.6pt},
    error bars/error mark options={tcRule, rotate=90, mark size=2pt}},
}


%----------------------------------------------------------------------
% 4. TABLE STYLE
%----------------------------------------------------------------------
% booktabs rules only — no vertical lines, no \hline.
% Row height and column padding are set once here; individual tables
% must not override them.
%----------------------------------------------------------------------
\renewcommand{\arraystretch}{1.15}
\setlength{\tabcolsep}{6pt}

\setlength{\heavyrulewidth}{0.9pt}
\setlength{\lightrulewidth}{0.5pt}
\setlength{\cmidrulewidth}{0.4pt}
\setlength{\aboverulesep}{0.35ex}
\setlength{\belowrulesep}{0.55ex}

% Every table body is set one step below body text. Individual tables must
% not declare their own size — this keeps the type size identical throughout.
\usepackage{etoolbox}
\AtBeginEnvironment{table}{\small}

% For a table too wide to fit \textwidth at \small. Use this instead of
% \resizebox, which scales the font by an arbitrary factor and therefore
% makes every wide table a different size.
\newcommand{\tabwide}{\footnotesize\setlength{\tabcolsep}{4pt}}

% Header cell. Use for every column heading so bolding stays consistent.
\newcommand{\thd}[1]{\textbf{#1}}

% Ragged-right paragraph column: L{3cm} instead of p{3cm}. Narrow justified
% columns produce heavy hyphenation and wide word gaps; ragged right does not.
\newcolumntype{L}[1]{>{\raggedright\arraybackslash}p{#1}}

% Standard note block underneath a table.
\newcommand{\tabnote}[1]{%
  \vspace{1.5mm}%
  \begin{minipage}{\textwidth}\footnotesize\textit{Note.} #1\end{minipage}}

% (The centred tabularx column type C is already defined in head.tex.)


%----------------------------------------------------------------------
% 5. CAPTIONS
%----------------------------------------------------------------------
\captionsetup{font=small, labelfont=bf, justification=justified,
              singlelinecheck=true}
\captionsetup[figure]{skip=8pt}
\captionsetup[table]{skip=6pt}

%======================================================================


%----------------------------------------------------------------------
% 1. COLOUR PALETTE
%----------------------------------------------------------------------
% Okabe-Ito palette: distinguishable under all common forms of colour
% vision deficiency. The L* value in each comment is the CIE lightness;
% series that must stay apart in greyscale printing are paired so that
% their L* values differ by at least 20.
%----------------------------------------------------------------------
\definecolor{tcBlue}{HTML}{0072B2}        % L* 47  — primary series
\definecolor{tcVermillion}{HTML}{D55E00}  % L* 52  — secondary series
\definecolor{tcTeal}{HTML}{009E73}        % L* 58  — third series
\definecolor{tcAmber}{HTML}{E69F00}       % L* 71  — second bar series
\definecolor{tcSky}{HTML}{56B4E9}         % L* 72  — spare
\definecolor{tcPurple}{HTML}{CC79A7}      % L* 61  — spare
\definecolor{tcGrey}{HTML}{595959}        % L* 38  — literature / prior work
\definecolor{tcGridGrey}{HTML}{BFBFBF}    % grid lines
\definecolor{tcRule}{HTML}{404040}        % axis lines and table rules

% Institutional colour, kept for the title page and header only.
% (HTWGreen is defined in style.tex and is deliberately not used as a
% data colour: it is hard to separate from teal in greyscale.)


%----------------------------------------------------------------------
% 2. SEMANTIC SERIES STYLES
%----------------------------------------------------------------------
% Colour carries MEANING and is constant across the whole thesis:
%
%   primary    blue        this work, main condition   (French, 0.5 mL,
%                                                       as-analyzed, pristine)
%   secondary  vermillion  this work, second condition (Algerian, 0.1 mL)
%   tertiary   teal        this work, third condition
%   emphasis   amber       derived / corrected quantity in BAR charts
%                          (ash-free, activated, sulfur-loaded)
%   reference  grey        literature or the preceding project work
%
% Shape and dash pattern duplicate that information so the figures stay
% readable in black-and-white print.
%----------------------------------------------------------------------
\pgfplotsset{
  % --- line series -----------------------------------------------------
  s primary/.style={
    tcBlue, thick, mark=*, mark size=2.3pt,
    mark options={fill=tcBlue, draw=tcBlue}},
  s secondary/.style={
    tcVermillion, thick, mark=square*, mark size=2.2pt,
    mark options={fill=tcVermillion, draw=tcVermillion}},
  s tertiary/.style={
    tcTeal, thick, mark=triangle*, mark size=2.8pt,
    mark options={fill=tcTeal, draw=tcTeal}},
  s quaternary/.style={
    tcAmber, thick, mark=diamond*, mark size=2.8pt,
    mark options={fill=tcAmber, draw=tcAmber}},
  s reference/.style={
    tcGrey, thick, densely dashed, mark=square*, mark size=2.2pt,
    mark options={solid, fill=tcGrey, draw=tcGrey}},
  % Open marker + dashed line = the same quantity on a different basis
  % (e.g. wet vs. moisture-corrected). Hue stays with the condition.
  s open/.style={
    dashed, mark options={solid, fill=white, line width=0.8pt}},
  %
  % --- bar series ------------------------------------------------------
  b primary/.style={fill=tcBlue,      draw=tcBlue!70!black,   line width=0.4pt},
  b emphasis/.style={fill=tcAmber,    draw=tcAmber!70!black,  line width=0.4pt},
  b secondary/.style={fill=tcVermillion, draw=tcVermillion!70!black, line width=0.4pt},
  b reference/.style={fill=tcGrey!30, draw=tcGrey,            line width=0.4pt},
}


%----------------------------------------------------------------------
% 3. AXIS STYLES
%----------------------------------------------------------------------
% One geometry for every plot in the thesis. The in-axis title is set once
% here, so every figure title is the same size and sits the same distance
% above the axis.
%----------------------------------------------------------------------
\pgfplotsset{
  thesisaxis/.style={
    width=0.86\textwidth,
    height=0.54\textwidth,
    title style={font=\normalsize, yshift=2pt},
    axis lines*=box,
    axis line style={tcRule, line width=0.5pt},
    tick align=outside,
    major tick length=3pt,
    minor tick num=0,
    tick style={tcRule, line width=0.4pt},
    ymajorgrids=true,
    grid style={tcGridGrey, line width=0.3pt},
    tick label style={font=\small},
    label style={font=\small},
    legend cell align=left,
    legend columns=2,
    legend style={
      font=\small,
      draw=tcGridGrey,
      fill=white,
      at={(0.5,-0.20)}, anchor=north,
      /tikz/every even column/.append style={column sep=0.6cm}},
    error bars/error bar style={tcRule, line width=0.6pt},
    error bars/error mark options={tcRule, rotate=90, mark size=2pt},
  },
  % line / scatter plots
  thesisline/.style={thesisaxis},
  % square-ish panel for scatter diagrams (e.g. van Krevelen)
  thesissquare/.style={
    thesisaxis,
    width=0.66\textwidth, height=0.54\textwidth,
    grid=both,
    minor grid style={tcGridGrey!50, line width=0.2pt}},
  % grouped bar charts
  thesisbar/.style={
    thesisaxis,
    ybar,
    area legend,
    bar width=18pt,
    ymin=0,
    enlarge x limits=0.28,
    xtick=data,
    tick pos=left,
    clip=false},
  % value labels on top of bars
  barlabels/.style={
    nodes near coords,
    nodes near coords align={vertical},
    every node near coord/.append style={
      font=\footnotesize, text=black,
      /pgf/number format/.cd, fixed, fixed zerofill, precision=1}},
  % overlay axis for a second y scale (no grid, no x axis)
  thesisaxis right/.style={
    width=0.86\textwidth,
    height=0.54\textwidth,
    axis y line*=right,
    axis x line=none,
    axis line style={tcRule, line width=0.5pt},
    tick align=outside,
    major tick length=3pt,
    tick style={tcRule, line width=0.4pt},
    tick label style={font=\small},
    label style={font=\small},
    error bars/error bar style={tcRule, line width=0.6pt},
    error bars/error mark options={tcRule, rotate=90, mark size=2pt}},
}


%----------------------------------------------------------------------
% 4. TABLE STYLE
%----------------------------------------------------------------------
% booktabs rules only — no vertical lines, no \hline.
% Row height and column padding are set once here; individual tables
% must not override them.
%----------------------------------------------------------------------
\renewcommand{\arraystretch}{1.15}
\setlength{\tabcolsep}{6pt}

\setlength{\heavyrulewidth}{0.9pt}
\setlength{\lightrulewidth}{0.5pt}
\setlength{\cmidrulewidth}{0.4pt}
\setlength{\aboverulesep}{0.35ex}
\setlength{\belowrulesep}{0.55ex}

% Every table body is set one step below body text. Individual tables must
% not declare their own size — this keeps the type size identical throughout.
\usepackage{etoolbox}
\AtBeginEnvironment{table}{\small}

% For a table too wide to fit \textwidth at \small. Use this instead of
% \resizebox, which scales the font by an arbitrary factor and therefore
% makes every wide table a different size.
\newcommand{\tabwide}{\footnotesize\setlength{\tabcolsep}{4pt}}

% Header cell. Use for every column heading so bolding stays consistent.
\newcommand{\thd}[1]{\textbf{#1}}

% Ragged-right paragraph column: L{3cm} instead of p{3cm}. Narrow justified
% columns produce heavy hyphenation and wide word gaps; ragged right does not.
\newcolumntype{L}[1]{>{\raggedright\arraybackslash}p{#1}}

% Standard note block underneath a table.
\newcommand{\tabnote}[1]{%
  \vspace{1.5mm}%
  \begin{minipage}{\textwidth}\footnotesize\textit{Note.} #1\end{minipage}}

% (The centred tabularx column type C is already defined in head.tex.)


%----------------------------------------------------------------------
% 5. CAPTIONS
%----------------------------------------------------------------------
\captionsetup{font=small, labelfont=bf, justification=justified,
              singlelinecheck=true}
\captionsetup[figure]{skip=8pt}
\captionsetup[table]{skip=6pt}

%======================================================================


%----------------------------------------------------------------------
% 1. COLOUR PALETTE
%----------------------------------------------------------------------
% Okabe-Ito palette: distinguishable under all common forms of colour
% vision deficiency. The L* value in each comment is the CIE lightness;
% series that must stay apart in greyscale printing are paired so that
% their L* values differ by at least 20.
%----------------------------------------------------------------------
\definecolor{tcBlue}{HTML}{0072B2}        % L* 47  — primary series
\definecolor{tcVermillion}{HTML}{D55E00}  % L* 52  — secondary series
\definecolor{tcTeal}{HTML}{009E73}        % L* 58  — third series
\definecolor{tcAmber}{HTML}{E69F00}       % L* 71  — second bar series
\definecolor{tcSky}{HTML}{56B4E9}         % L* 72  — spare
\definecolor{tcPurple}{HTML}{CC79A7}      % L* 61  — spare
\definecolor{tcGrey}{HTML}{595959}        % L* 38  — literature / prior work
\definecolor{tcGridGrey}{HTML}{BFBFBF}    % grid lines
\definecolor{tcRule}{HTML}{404040}        % axis lines and table rules

% Institutional colour, kept for the title page and header only.
% (HTWGreen is defined in style.tex and is deliberately not used as a
% data colour: it is hard to separate from teal in greyscale.)


%----------------------------------------------------------------------
% 2. SEMANTIC SERIES STYLES
%----------------------------------------------------------------------
% Colour carries MEANING and is constant across the whole thesis:
%
%   primary    blue        this work, main condition   (French, 0.5 mL,
%                                                       as-analyzed, pristine)
%   secondary  vermillion  this work, second condition (Algerian, 0.1 mL)
%   tertiary   teal        this work, third condition
%   emphasis   amber       derived / corrected quantity in BAR charts
%                          (ash-free, activated, sulfur-loaded)
%   reference  grey        literature or the preceding project work
%
% Shape and dash pattern duplicate that information so the figures stay
% readable in black-and-white print.
%----------------------------------------------------------------------
\pgfplotsset{
  % --- line series -----------------------------------------------------
  s primary/.style={
    tcBlue, thick, mark=*, mark size=2.3pt,
    mark options={fill=tcBlue, draw=tcBlue}},
  s secondary/.style={
    tcVermillion, thick, mark=square*, mark size=2.2pt,
    mark options={fill=tcVermillion, draw=tcVermillion}},
  s tertiary/.style={
    tcTeal, thick, mark=triangle*, mark size=2.8pt,
    mark options={fill=tcTeal, draw=tcTeal}},
  s quaternary/.style={
    tcAmber, thick, mark=diamond*, mark size=2.8pt,
    mark options={fill=tcAmber, draw=tcAmber}},
  s reference/.style={
    tcGrey, thick, densely dashed, mark=square*, mark size=2.2pt,
    mark options={solid, fill=tcGrey, draw=tcGrey}},
  % Open marker + dashed line = the same quantity on a different basis
  % (e.g. wet vs. moisture-corrected). Hue stays with the condition.
  s open/.style={
    dashed, mark options={solid, fill=white, line width=0.8pt}},
  %
  % --- bar series ------------------------------------------------------
  b primary/.style={fill=tcBlue,      draw=tcBlue!70!black,   line width=0.4pt},
  b emphasis/.style={fill=tcAmber,    draw=tcAmber!70!black,  line width=0.4pt},
  b secondary/.style={fill=tcVermillion, draw=tcVermillion!70!black, line width=0.4pt},
  b reference/.style={fill=tcGrey!30, draw=tcGrey,            line width=0.4pt},
}


%----------------------------------------------------------------------
% 3. AXIS STYLES
%----------------------------------------------------------------------
% One geometry for every plot in the thesis. The in-axis title is set once
% here, so every figure title is the same size and sits the same distance
% above the axis.
%----------------------------------------------------------------------
\pgfplotsset{
  thesisaxis/.style={
    width=0.86\textwidth,
    height=0.54\textwidth,
    title style={font=\normalsize, yshift=2pt},
    axis lines*=box,
    axis line style={tcRule, line width=0.5pt},
    tick align=outside,
    major tick length=3pt,
    minor tick num=0,
    tick style={tcRule, line width=0.4pt},
    ymajorgrids=true,
    grid style={tcGridGrey, line width=0.3pt},
    tick label style={font=\small},
    label style={font=\small},
    legend cell align=left,
    legend columns=2,
    legend style={
      font=\small,
      draw=tcGridGrey,
      fill=white,
      at={(0.5,-0.20)}, anchor=north,
      /tikz/every even column/.append style={column sep=0.6cm}},
    error bars/error bar style={tcRule, line width=0.6pt},
    error bars/error mark options={tcRule, rotate=90, mark size=2pt},
  },
  % line / scatter plots
  thesisline/.style={thesisaxis},
  % square-ish panel for scatter diagrams (e.g. van Krevelen)
  thesissquare/.style={
    thesisaxis,
    width=0.66\textwidth, height=0.54\textwidth,
    grid=both,
    minor grid style={tcGridGrey!50, line width=0.2pt}},
  % grouped bar charts
  thesisbar/.style={
    thesisaxis,
    ybar,
    area legend,
    bar width=18pt,
    ymin=0,
    enlarge x limits=0.28,
    xtick=data,
    tick pos=left,
    clip=false},
  % value labels on top of bars
  barlabels/.style={
    nodes near coords,
    nodes near coords align={vertical},
    every node near coord/.append style={
      font=\footnotesize, text=black,
      /pgf/number format/.cd, fixed, fixed zerofill, precision=1}},
  % overlay axis for a second y scale (no grid, no x axis)
  thesisaxis right/.style={
    width=0.86\textwidth,
    height=0.54\textwidth,
    axis y line*=right,
    axis x line=none,
    axis line style={tcRule, line width=0.5pt},
    tick align=outside,
    major tick length=3pt,
    tick style={tcRule, line width=0.4pt},
    tick label style={font=\small},
    label style={font=\small},
    error bars/error bar style={tcRule, line width=0.6pt},
    error bars/error mark options={tcRule, rotate=90, mark size=2pt}},
}


%----------------------------------------------------------------------
% 4. TABLE STYLE
%----------------------------------------------------------------------
% booktabs rules only — no vertical lines, no \hline.
% Row height and column padding are set once here; individual tables
% must not override them.
%----------------------------------------------------------------------
\renewcommand{\arraystretch}{1.15}
\setlength{\tabcolsep}{6pt}

\setlength{\heavyrulewidth}{0.9pt}
\setlength{\lightrulewidth}{0.5pt}
\setlength{\cmidrulewidth}{0.4pt}
\setlength{\aboverulesep}{0.35ex}
\setlength{\belowrulesep}{0.55ex}

% Every table body is set one step below body text. Individual tables must
% not declare their own size — this keeps the type size identical throughout.
\usepackage{etoolbox}
\AtBeginEnvironment{table}{\small}

% For a table too wide to fit \textwidth at \small. Use this instead of
% \resizebox, which scales the font by an arbitrary factor and therefore
% makes every wide table a different size.
\newcommand{\tabwide}{\footnotesize\setlength{\tabcolsep}{4pt}}

% Header cell. Use for every column heading so bolding stays consistent.
\newcommand{\thd}[1]{\textbf{#1}}

% Ragged-right paragraph column: L{3cm} instead of p{3cm}. Narrow justified
% columns produce heavy hyphenation and wide word gaps; ragged right does not.
\newcolumntype{L}[1]{>{\raggedright\arraybackslash}p{#1}}

% Standard note block underneath a table.
\newcommand{\tabnote}[1]{%
  \vspace{1.5mm}%
  \begin{minipage}{\textwidth}\footnotesize\textit{Note.} #1\end{minipage}}

% (The centred tabularx column type C is already defined in head.tex.)


%----------------------------------------------------------------------
% 5. CAPTIONS
%----------------------------------------------------------------------
\captionsetup{font=small, labelfont=bf, justification=justified,
              singlelinecheck=true}
\captionsetup[figure]{skip=8pt}
\captionsetup[table]{skip=6pt}

%======================================================================


%----------------------------------------------------------------------
% 1. COLOUR PALETTE
%----------------------------------------------------------------------
% Okabe-Ito palette: distinguishable under all common forms of colour
% vision deficiency. The L* value in each comment is the CIE lightness;
% series that must stay apart in greyscale printing are paired so that
% their L* values differ by at least 20.
%----------------------------------------------------------------------
\definecolor{tcBlue}{HTML}{0072B2}        % L* 47  — primary series
\definecolor{tcVermillion}{HTML}{D55E00}  % L* 52  — secondary series
\definecolor{tcTeal}{HTML}{009E73}        % L* 58  — third series
\definecolor{tcAmber}{HTML}{E69F00}       % L* 71  — second bar series
\definecolor{tcSky}{HTML}{56B4E9}         % L* 72  — spare
\definecolor{tcPurple}{HTML}{CC79A7}      % L* 61  — spare
\definecolor{tcGrey}{HTML}{595959}        % L* 38  — literature / prior work
\definecolor{tcGridGrey}{HTML}{BFBFBF}    % grid lines
\definecolor{tcRule}{HTML}{404040}        % axis lines and table rules

% Institutional colour, kept for the title page and header only.
% (HTWGreen is defined in style.tex and is deliberately not used as a
% data colour: it is hard to separate from teal in greyscale.)


%----------------------------------------------------------------------
% 2. SEMANTIC SERIES STYLES
%----------------------------------------------------------------------
% Colour carries MEANING and is constant across the whole thesis:
%
%   primary    blue        this work, main condition   (French, 0.5 mL,
%                                                       as-analyzed, pristine)
%   secondary  vermillion  this work, second condition (Algerian, 0.1 mL)
%   tertiary   teal        this work, third condition
%   emphasis   amber       derived / corrected quantity in BAR charts
%                          (ash-free, activated, sulfur-loaded)
%   reference  grey        literature or the preceding project work
%
% Shape and dash pattern duplicate that information so the figures stay
% readable in black-and-white print.
%----------------------------------------------------------------------
\pgfplotsset{
  % --- line series -----------------------------------------------------
  s primary/.style={
    tcBlue, thick, mark=*, mark size=2.3pt,
    mark options={fill=tcBlue, draw=tcBlue}},
  s secondary/.style={
    tcVermillion, thick, mark=square*, mark size=2.2pt,
    mark options={fill=tcVermillion, draw=tcVermillion}},
  s tertiary/.style={
    tcTeal, thick, mark=triangle*, mark size=2.8pt,
    mark options={fill=tcTeal, draw=tcTeal}},
  s quaternary/.style={
    tcAmber, thick, mark=diamond*, mark size=2.8pt,
    mark options={fill=tcAmber, draw=tcAmber}},
  s reference/.style={
    tcGrey, thick, densely dashed, mark=square*, mark size=2.2pt,
    mark options={solid, fill=tcGrey, draw=tcGrey}},
  % Open marker + dashed line = the same quantity on a different basis
  % (e.g. wet vs. moisture-corrected). Hue stays with the condition.
  s open/.style={
    dashed, mark options={solid, fill=white, line width=0.8pt}},
  %
  % --- bar series ------------------------------------------------------
  b primary/.style={fill=tcBlue,      draw=tcBlue!70!black,   line width=0.4pt},
  b emphasis/.style={fill=tcAmber,    draw=tcAmber!70!black,  line width=0.4pt},
  b secondary/.style={fill=tcVermillion, draw=tcVermillion!70!black, line width=0.4pt},
  b reference/.style={fill=tcGrey!30, draw=tcGrey,            line width=0.4pt},
}


%----------------------------------------------------------------------
% 3. AXIS STYLES
%----------------------------------------------------------------------
% One geometry for every plot in the thesis. The in-axis title is set once
% here, so every figure title is the same size and sits the same distance
% above the axis.
%----------------------------------------------------------------------
\pgfplotsset{
  thesisaxis/.style={
    width=0.86\textwidth,
    height=0.54\textwidth,
    title style={font=\normalsize, yshift=2pt},
    axis lines*=box,
    axis line style={tcRule, line width=0.5pt},
    tick align=outside,
    major tick length=3pt,
    minor tick num=0,
    tick style={tcRule, line width=0.4pt},
    ymajorgrids=true,
    grid style={tcGridGrey, line width=0.3pt},
    tick label style={font=\small},
    label style={font=\small},
    legend cell align=left,
    legend columns=2,
    legend style={
      font=\small,
      draw=tcGridGrey,
      fill=white,
      at={(0.5,-0.20)}, anchor=north,
      /tikz/every even column/.append style={column sep=0.6cm}},
    error bars/error bar style={tcRule, line width=0.6pt},
    error bars/error mark options={tcRule, rotate=90, mark size=2pt},
  },
  % line / scatter plots
  thesisline/.style={thesisaxis},
  % square-ish panel for scatter diagrams (e.g. van Krevelen)
  thesissquare/.style={
    thesisaxis,
    width=0.66\textwidth, height=0.54\textwidth,
    grid=both,
    minor grid style={tcGridGrey!50, line width=0.2pt}},
  % grouped bar charts
  thesisbar/.style={
    thesisaxis,
    ybar,
    area legend,
    bar width=18pt,
    ymin=0,
    enlarge x limits=0.28,
    xtick=data,
    tick pos=left,
    clip=false},
  % value labels on top of bars
  barlabels/.style={
    nodes near coords,
    nodes near coords align={vertical},
    every node near coord/.append style={
      font=\footnotesize, text=black,
      /pgf/number format/.cd, fixed, fixed zerofill, precision=1}},
  % overlay axis for a second y scale (no grid, no x axis)
  thesisaxis right/.style={
    width=0.86\textwidth,
    height=0.54\textwidth,
    axis y line*=right,
    axis x line=none,
    axis line style={tcRule, line width=0.5pt},
    tick align=outside,
    major tick length=3pt,
    tick style={tcRule, line width=0.4pt},
    tick label style={font=\small},
    label style={font=\small},
    error bars/error bar style={tcRule, line width=0.6pt},
    error bars/error mark options={tcRule, rotate=90, mark size=2pt}},
}


%----------------------------------------------------------------------
% 4. TABLE STYLE
%----------------------------------------------------------------------
% booktabs rules only — no vertical lines, no \hline.
% Row height and column padding are set once here; individual tables
% must not override them.
%----------------------------------------------------------------------
\renewcommand{\arraystretch}{1.15}
\setlength{\tabcolsep}{6pt}

\setlength{\heavyrulewidth}{0.9pt}
\setlength{\lightrulewidth}{0.5pt}
\setlength{\cmidrulewidth}{0.4pt}
\setlength{\aboverulesep}{0.35ex}
\setlength{\belowrulesep}{0.55ex}

% Every table body is set one step below body text. Individual tables must
% not declare their own size — this keeps the type size identical throughout.
\usepackage{etoolbox}
\AtBeginEnvironment{table}{\small}

% For a table too wide to fit \textwidth at \small. Use this instead of
% \resizebox, which scales the font by an arbitrary factor and therefore
% makes every wide table a different size.
\newcommand{\tabwide}{\footnotesize\setlength{\tabcolsep}{4pt}}

% Header cell. Use for every column heading so bolding stays consistent.
\newcommand{\thd}[1]{\textbf{#1}}

% Ragged-right paragraph column: L{3cm} instead of p{3cm}. Narrow justified
% columns produce heavy hyphenation and wide word gaps; ragged right does not.
\newcolumntype{L}[1]{>{\raggedright\arraybackslash}p{#1}}

% Standard note block underneath a table.
\newcommand{\tabnote}[1]{%
  \vspace{1.5mm}%
  \begin{minipage}{\textwidth}\footnotesize\textit{Note.} #1\end{minipage}}

% (The centred tabularx column type C is already defined in head.tex.)


%----------------------------------------------------------------------
% 5. CAPTIONS
%----------------------------------------------------------------------
\captionsetup{font=small, labelfont=bf, justification=justified,
              singlelinecheck=true}
\captionsetup[figure]{skip=8pt}
\captionsetup[table]{skip=6pt}
